\chapter{Numerical verification and scope}\label{ch-numerical}

This chapter reports computations. They play no role in any proof: every
theorem in this thesis is proved from the axioms, and the experiments are
recorded because they locate the extremal graphs, because they quantify how
much slack the bound has at small orders, and because they delimit the class
of matrices to which the bound extends. Each table is labelled
\emph{numerical} in the sense of Section~\ref{sec-notation}, except where an
exact identity is stated and proved alongside it.

\section{Protocol}\label{sec-protocol}
Three kinds of computation were run, with deliberately separated roles.

\emph{Exact arithmetic.} Every displayed identity, factorisation, completion
of squares and positivity assertion in Chapters~\ref{ch-moment}
and~\ref{ch-bound} was re-checked with exact rationals and with elements of
the quadratic fields $\mathbb Q(\sqrt3)$ and $\mathbb Q(\sqrt5)$ represented
as pairs of rationals, never in floating point. The constants $c$, $\cstar$,
$(3+\sqrt5)/4$ and $3/\sqrt5$ were recomputed to sixty decimal places:
\begin{align*}
 c&=1.284673255723918505060732785864589260927238391290291280182\ldots,\\
 \cstar&=1.281568196872167876752435013991917240681756712288900353890\ldots,\\
 (3+\sqrt5)/4&=1.309016994374947424102293417182819058860154589902881431067\ldots,\\
 3/\sqrt5&=1.341640786499873817845504201238765741264371015766915434562\ldots,
\end{align*}
from which $c-\cstar=0.0031050588517506283082977718726720202454816\ldots$ and
$(3+\sqrt5)/4-c=0.024343738651028919041560631318229797932916\ldots$. The
six-part quotient value at the optimising weight was confirmed equal to
$\cstar$ to forty digits.

\emph{Exhaustive enumeration.} For $3\le n\le7$ all $2^{n(n-1)/2}$ labelled
graphs were enumerated, $\Lambda_3$ was computed with a symmetric eigensolver
in double precision, and the result was compared with $cn-3$ evaluated from
the exact expression for $c$. A $7\times7$ symmetric eigenvalue computation in
double precision has backward error of order $10^{-15}$, while the smallest
margin encountered is $2.2\times10^{-2}$; the conclusions are therefore
insensitive to rounding.

\emph{Randomised sampling.} For the weighted tests, symmetric matrices with
zero diagonal and independent uniform $[0,1]$ upper-triangular entries were
sampled with a fixed seed, so the runs are reproducible.

\section{Exhaustive verification for $n\le7$}\label{sec-exhaustive}
The total number of labelled graphs examined is
$8+64+1024+32768+2097152=2\,131\,016$. There were no violations. The extremal
values, the extremal graphs up to isomorphism and the number of labelled
copies are in Table~\ref{tab-exhaustive}. At $n=3$ all eight graphs tie, since $\Lambda_3=\tr A_G=0$ there;
at $n=4$ the extremal graph is unique; at $n=5$ there are exactly two
extremal isomorphism classes; at $n=6$ and $n=7$ the extremal graph is again
unique.

\begin{table}[htbp]
\centering
\singlespacing
\caption{Exhaustive enumeration over all labelled graphs, $3\le n\le7$.
``Margin'' is $cn-3-\max\Lambda_3$; the exact value of $\max\Lambda_3$ is
printed where it is known in radicals. At $n=3$ every graph satisfies
$\Lambda_3=\tr A_G=0$.}\label{tab-exhaustive}
\footnotesize
\begin{tabular}{@{}r r l r r l@{}}
\toprule
$n$ & \#\ graphs & $\max\Lambda_3$ & $cn-3$ & margin & extremal graph \\
\midrule
3 & 8 & $0$ & 0.854019767 & 0.854019767 & --- \\
4 & 64 & $2$ & 2.138693023 & 0.138693023 & $C_4$ \\
5 & 1024 & $1+\sqrt5=3.236067977$ & 3.423366279 & 0.187298301 & $C_5$, $K_1\vee C_4$ \\
6 & 32768 & $\sqrt5+\sqrt6=4.685557720$ & 4.708039534 & 0.022481814 & $K_1\vee C_5$ \\
7 & 2097152 & $\tfrac{1+\sqrt{41}}2+\sqrt5=5.937630096$ & 5.992712790 & 0.055082694 & $K_2\vee C_5$ \\
\bottomrule
\end{tabular}
\end{table}

The number of labelled maximisers is $3$ at $n=4$, $27$ at $n=5$, $72$ at
$n=6$ and $252$ at $n=7$. The value $72=6!/10$ matches
$|\operatorname{Aut}(K_1\vee C_5)|=10$, and $252=7!/20$ matches
$|\operatorname{Aut}(K_2\vee C_5)|=2\cdot10=20$; for $n=7$ all $252$
maximisers were additionally checked to have degree sequence
$(6,6,4,4,4,4,4)$ and to induce a $5$-cycle on the five vertices of degree
four. The split $27=12+15$ at $n=5$ matches $5!/10=12$ labelled copies of
$C_5$ and $15$ labelled copies of $K_1\vee C_4$; the two classes are not
isomorphic but are exactly tied, which is why uniqueness fails there.

Two features of Table~\ref{tab-exhaustive} deserve comment. First, the order
$n=3$ is the smallest for which the bound is even meaningful: for $n=1,2$ the
right-hand side $cn-3$ is negative while $\Lambda_3=\tr A_G=0$ whenever
$n\le3$, so the hypothesis $n\ge3$ in Theorem~\ref{main} cannot be dropped.
Second, the margin is smallest at $n=6$, which is exactly the order of the
template $K_1\vee C_5$ that produces $\cstar$; the bound is therefore closest
to being tight precisely where the lower-bound construction lives, which is
the numerical counterpart of the gap $c-\cstar$.

\section{Blow-up families}\label{sec-blowups}
The template $H=K_1\vee C_5$ generates an explicit infinite family whose
$\Lambda_3$ can be computed in closed form.

\begin{proposition}[Equal blow-ups of $K_1\vee C_5$]\label{equal-blowup}
Let $q\ge1$ and let $G=H[q]$ be the graph on $n=6q$ vertices obtained by
replacing every vertex of $H$ by a clique of size $q$ and every edge of $H$ by
a complete bipartite graph. Then
\[
 \Lambda_3(G)=\frac{\sqrt5+\sqrt6+3}{6}\,n-3,
 \qquad
 \frac{\sqrt5+\sqrt6+3}{6}=1.2809262867\ldots .
\]
\end{proposition}
\begin{proof}
By Lemma~\ref{blowup-spec} the spectrum of $G$ consists of $q\lambda+q-1$ for
$\lambda\in\operatorname{spec}H$ together with $-1$ repeated $6(q-1)$ times.
The spectrum of $H$ is
$1+\sqrt6$, $\tfrac{\sqrt5-1}2$ twice, $1-\sqrt6$ and
$-\tfrac{1+\sqrt5}2$ twice, so the three largest eigenvalues of $G$ are
$q(1+\sqrt6)+q-1$ and $q\tfrac{\sqrt5-1}2+q-1$ twice; for $q\ge1$ these
exceed both $q(1-\sqrt6)+q-1$ and $-1$. Their sum is
$q(\sqrt6+\sqrt5+3)-3$, which equals the claimed expression at $n=6q$.
\end{proof}

The coefficient $1.2809262867\ldots$ is strictly below both $\cstar$ and $c$.
Equal part sizes are therefore suboptimal, and the loss is quantified:
\[
 \cstar-\frac{\sqrt5+\sqrt6+3}{6}=0.0006419102\ldots,
 \qquad
 c-\frac{\sqrt5+\sqrt6+3}{6}=0.0037469690\ldots .
\]
Table~\ref{tab-blowup} confirms the closed form numerically and shows the
distance to both comparison lines.

\begin{table}[htbp]
\centering
\singlespacing
\caption{Equal blow-ups $G=H[q]$ of $H=K_1\vee C_5$, $n=6q$. All values
numerical; the column $\Lambda_3$ agrees with the closed form of
Proposition~\ref{equal-blowup} to the printed precision.}\label{tab-blowup}
\small
\begin{tabular}{@{}r r r r r r@{}}
\toprule
$q$ & $n$ & $\Lambda_3(G)$ & $\Lambda_3/n$ & $\Lambda_3-(cn-3)$ & $\Lambda_3-(\cstar n-3)$ \\
\midrule
1 & 6 & 4.685558 & 0.780926 & $-0.022482$ & $-0.003851$ \\
2 & 12 & 12.371115 & 1.030926 & $-0.044964$ & $-0.007703$ \\
3 & 18 & 20.056673 & 1.114260 & $-0.067445$ & $-0.011554$ \\
5 & 30 & 35.427789 & 1.180926 & $-0.112409$ & $-0.019257$ \\
10 & 60 & 73.855577 & 1.230926 & $-0.224818$ & $-0.038515$ \\
20 & 120 & 150.711154 & 1.255926 & $-0.449636$ & $-0.077029$ \\
50 & 300 & 381.277886 & 1.270926 & $-1.124091$ & $-0.192573$ \\
100 & 600 & 765.555772 & 1.275926 & $-2.248181$ & $-0.385146$ \\
250 & 1500 & 1918.389430 & 1.278926 & $-5.620454$ & $-0.962865$ \\
500 & 3000 & 3839.778860 & 1.279926 & $-11.240907$ & $-1.925730$ \\
\bottomrule
\end{tabular}
\end{table}

The optimal part sizes are given by Lemma~\ref{attainment}: apex mass
$p^*=(91-25\sqrt7)/126=0.1972715653\ldots$ and rim mass
$b^*=(7+5\sqrt7)/126=0.1605456869\ldots$ for each of the five rim parts.
Realising these as integers introduces a rounding error of order one vertex per
part, so the resulting graphs fall short of $\cstar n-3$ by an amount that
decreases with $n$. Table~\ref{tab-optblowup} records this convergence, using
the general quotient form~\eqref{quotient-unequal}.

\begin{table}[htbp]
\centering
\singlespacing
\caption{Optimally weighted blow-ups of $K_1\vee C_5$: one apex clique and
five rim cliques, sizes rounded from $(p^*,b^*,\ldots,b^*)n$. The last column
tends to $0$, which is the numerical form of $\Sigma_3(S)=\cstar$ for the
optimal six-part step graphon $S$.}\label{tab-optblowup}
\small
\begin{tabular}{@{}r l r r r r@{}}
\toprule
$n$ & part sizes & $\Lambda_3$ & $\Lambda_3/n$ & $\Lambda_3-(cn-3)$ & $\Lambda_3-(\cstar n-3)$ \\
\midrule
600 & $(120,96^5)$ & 765.937967 & 1.276563 & $-1.865986$ & $-0.002951$ \\
1200 & $(235,193^5)$ & 1534.880189 & 1.279067 & $-3.727718$ & $-0.001647$ \\
3000 & $(590,482^5)$ & 3841.703863 & 1.280568 & $-9.315904$ & $-0.000728$ \\
\bottomrule
\end{tabular}
\end{table}

The two tables together exhibit the structure of the gap. The optimally
weighted family approaches $\cstar n-3$ from below and is far from $cn-3$,
while the bound $cn-3$ is not attained by any graph found here. Nothing in
this chapter decides whether a different family could exceed $\cstar n-3$; by
Proposition~\ref{no-gain} no family of weighted graphs can exceed
$\mathcal M n-3$ asymptotically, and $\mathcal M$ itself is only bracketed.

\section{Weighted matrices}\label{sec-num-weighted}
Proposition~\ref{weighted-form} extends the bound to every symmetric matrix
with zero diagonal and off-diagonal entries in $[0,1]$. Table~\ref{tab-weighted}
reports randomised tests of that statement. No violation occurred. A control run of $200$ matrices of order $30$ with entries uniform in
$[0,2]$ produced $198$ violations of $\Lambda_3\le cn-3$, so the range
restriction $A_{ij}\le1$ is the operative hypothesis and not an artefact of
the sampling; this is the numerical counterpart of Remark~\ref{weighted-scope}.

\begin{table}[htbp]
\centering
\singlespacing
\caption{Random $[0,1]$-weighted loopless matrices, $4000$ samples per order,
fixed seed. ``Max'' is the largest $\Lambda_3$ observed; ``bound'' is $cn-3$.
No violation occurred at any order.}\label{tab-weighted}
\small
\begin{tabular}{@{}r r r r r@{}}
\toprule
$n$ & samples & violations & max $\Lambda_3$ & $cn-3$ \\
\midrule
4 & 4000 & 0 & 1.663987 & 2.138693 \\
6 & 4000 & 0 & 3.637893 & 4.708040 \\
9 & 4000 & 0 & 6.244673 & 8.562059 \\
15 & 4000 & 0 & 10.708949 & 16.270099 \\
30 & 4000 & 0 & 20.579009 & 35.540198 \\
60 & 4000 & 0 & 38.337336 & 74.080395 \\
\bottomrule
\end{tabular}
\end{table}

The random weighted matrices are far from extremal: at $n=60$ the observed
maximum is barely half the bound. This is expected, because a random weighted
matrix has a dense, near-semicircular spectrum whose third eigenvalue is small
compared with $n$, whereas the extremal constructions are sparse-template
blow-ups with three large eigenvalues. The experiment is therefore a test of
the hypothesis range, not a search for extremals.

\section{What the computations do not show}\label{sec-num-limits}
Three limitations are stated explicitly so that no reader mistakes numerical
evidence for a theorem.

\begin{itemize}[leftmargin=1.6em,itemsep=4pt]
\item Exhaustive verification up to $n=7$ says nothing about $n\ge8$; the
margin at $n=6$ is small enough that no extrapolation from the table is
possible.
\item The absence of violations in Table~\ref{tab-weighted} does not prove
Proposition~\ref{weighted-form}; that proposition is proved in
Chapter~\ref{ch-bound} by Lemma~\ref{bridge}, and the sampling only checks
that the proof's hypotheses were not mis-stated.
\item The blow-up tables identify a candidate extremal family but do not
establish that $\mathcal M=\cstar$. Proposition~\ref{no-gain} shows that the
asymptotic constant is the same for weighted and unweighted problems; it does
not evaluate it.
\end{itemize}

\section{Reproducibility}\label{sec-repro}
The scripts that produced Tables~\ref{tab-exhaustive}
to~\ref{tab-weighted} and the exact-arithmetic checks of
Section~\ref{sec-protocol} are stored in the accompanying repository under
\texttt{thesis\_main/scripts/}, with the raw outputs as JSON in the same
directory. Each script is deterministic: the enumeration is exhaustive, and
the randomised tests use the fixed seed $20261104$. The exact-arithmetic
checks use only integer and quadratic-field operations and print their results
at sixty digits, so they can be reproduced without any numerical library.
